\section{A quantitative planar cycle-count consequence}
\label{sec:planar-cycle-error}

Aldred and Thomassen proved the planar asymptotic bound
$c(G)\le(1/2+o(1))2^{\beta(G)}$ in 2008.\footnote{R. E. L. Aldred and C. Thomassen,
\emph{On the maximum number of cycles in a planar graph}, Journal of Graph Theory 57 (2008), 255--264,
\url{https://doi.org/10.1002/jgt.20290}. Their institutional publication record is
\url{https://orbit.dtu.dk/en/publications/on-the-maximum-number-of-cycles-in-a-planar-graph/}.}
The argument below gives an explicit error term from the present machinery.
The original article's full proof was not available during this comparison.
We therefore make no priority claim for the quantitative rate.

\begin{proposition}[An explicit planar error]\label{prop:planar-cycle-error}
Every planar graph of rank $r\ge1$ satisfies
\[
 c(G)\le2^{r-1}+32257\frac{2^r}{r}.
\]
More generally, for graphs of orientable genus $g=o(r)$,
\[
 c(G)\le2^{r-1}+O(2^r/r),
\]
with an absolute implied constant for sufficiently large $r$.
\end{proposition}

\begin{lemma}[Embedded marked sets]\label{lem:embedded-marked-forest}
Let $H$ be a connected graph of rank $R$, embedded in an orientable surface of genus $g$.
Let $W$ be a nonempty union of cycles with $m$ edges.
For every integer $L\ge6$ such that
\[
 B=(L-5)R-(4L+12)m-2(L+1)g>0,
\]
one has
\[
 p(H,W)\le\frac12+
 \frac{3L(L+1)(L-1)2^{L-2}}{B}.
\]
\end{lemma}

\begin{proof}
Inside each vertex disk, replace a high-degree vertex by a trivalent tree respecting the external cyclic order.
Mark every new vertex incident with an edge of $W$.
Retain the original vertex when its degree is at most three, marking it if necessary.
The marked set $P$ has size $p\le2m$.
No connectivity condition on the marked set is needed.

Contracting the replacement trees identifies the two cycle spaces.
If the original even subgraph is a forest outside $W$, its lift is a forest after deleting $P$.
Indeed, a forest remains a forest when its vertices are replaced by trees.
Delete leaves and suppress degree-two vertices outside $P$.
These operations preserve the cycle space and the required forest implication.
They also preserve the embedding.
Call the resulting connected multigraph $\Gamma$.
Every unmarked vertex has degree three, while each marked vertex has degree at least one and at most three.
Consequently,
\[
 |V(\Gamma)|\le2R+2p-2,
 \qquad |E(\Gamma)|\le3R+2p-3.
\]

Put $J=\Gamma-P$.
Deleting $P$ removes at most $3p$ edges, giving $\beta(J)\ge R-2p$ when $J$ is nonempty.
When $J$ is empty, subcubicity gives $R\le p/2+1\le2p$, so the same bound holds.
The face argument of Lemma~\ref{lem:path-face-packing} now gives at least
\[
 R-2p-2g-\frac{6R+4p}{L+1}
 =\frac{(L-5)R-(2L+6)p-2(L+1)g}{L+1}
\]
retained faces with boundary length at most $L$.
Greedy selection yields $k$ vertex-disjoint cycles outside $P$, where
\[
 k\ge\frac{(L-5)R-(2L+6)p-2(L+1)g}{3L(L+1)}
 \ge\frac{B}{3L(L+1)}.
\]
Each cycle's vertex set has boundary at most $L$ in $\Gamma$.
Apply \eqref{eq:path-small-cuts} and the forest implication to obtain the stated bound.
\end{proof}

\begin{proof}[Proof of Proposition~\ref{prop:planar-cycle-error}]
Join different components by bridges, preserving the rank, cycles, and genus bound.
For $r\ge2$, put $k=\lceil\log_2r\rceil$.
Choose $k$ edges $D$ outside a spanning tree.
Exactly $2^{r-k}$ even subgraphs avoid $D$.

For each $e=uv\in D$, add a fresh vertex $x_e$ and the edges $ux_e,x_ev$.
Embed this new path beside $e$ in a narrow neighborhood of that edge.
The resulting graph $H$ has genus at most $g$ and rank $R=r+k$.
Let $W$ be the union of the $k$ resulting triangles, so $|E(W)|=3k$.

For each original cycle $C$ meeting $D$ and each subset of the new triangles, take their symmetric difference with $C$.
These $2^k$ even subgraphs have the same forest outside $W$.
The new edges determine the chosen subset and hence determine $C$.
Thus at most $p(H,W)2^r$ original cycles meet $D$.
The cycles avoiding $D$ contribute at most $2^{r-k}$.

Set $L=7$ in Lemma~\ref{lem:embedded-marked-forest}.
When the denominator is positive, it gives
\begin{equation}\label{eq:planar-error-denominator}
 \frac{c(G)}{2^r}
 \le\frac12+
 \frac{32256}{2r-118\lceil\log_2r\rceil-16g}
 +2^{-\lceil\log_2r\rceil}.
\end{equation}
If $g=0$ and $r\ge2048$, then $r\ge118\lceil\log_2r\rceil$.
The last display is therefore at most $1/2+32257/r$.
For $1\le r<2048$, the trivial bound $c(G)<2^r$ proves the same assertion.
If $g=o(r)$, the denominator in \eqref{eq:planar-error-denominator} is $(2-o(1))r$.
This proves the remaining assertion.
\end{proof}

The constants come from a fixed face-length cutoff and are not optimized.
This proof uses only embedding-preserving additions beside selected edges.

\section{A polynomial cycle-count error with a fixed excluded minor}
\label{sec:minor-cycle-error}

The same order of error holds throughout every fixed proper minor-closed class.
Here the separators are applied before the witness paths are added.
Thus the proof does not need to preserve minor exclusion under those additions.

\begin{theorem}[Cycle-count error with an excluded minor]\label{thm:minor-cycle-error}
For each fixed integer $h\ge4$, there is a constant $C_h$ such that every simple graph $G$
with no $K_h$ minor and rank $r\ge1$ satisfies
\[
 c(G)\le2^{r-1}+C_h\frac{2^r}{r}.
\]
One may take
\[
 K=4096h^3,\qquad C_h=1+8K(3K-1)2^{3K-2}.
\]
In particular, this bound holds in every fixed proper minor-closed class, with a class-dependent constant.
\end{theorem}

\begin{proof}
Suppose first that $c(G)>2^{r-1}$.
All cycles then lie in one block.
Indeed, two nonempty groups of cyclic blocks, with ranks $a,b\ge1$, would give
\[
 c(G)\le(2^a-1)+(2^b-1)\le2^{a+b-1}=2^{r-1}.
\]
Discard all other blocks.
This preserves the rank and cycle count, giving a two-connected graph $B$.

For each vertex $v$ of $B$, the graph $B-v$ is connected.
The incident pattern of a uniform even subgraph is therefore uniform over all even subsets of the incident edges.
A cycle uses zero or two edges at $v$.
If $v$ has degree $d$, this gives
\[
 \frac{c(B)}{2^r}\le\frac{1+\binom d2}{2^{d-1}}.
\]
The right side is at most one half when $d\ge6$.
Consequently, every vertex of $B$ has degree at most five.

It suffices to consider $r\ge512$.
Suppress all degree-two vertices of $B$, retaining parallel edges.
Since $r\ge2$, the resulting loopless multigraph $\Gamma$ has minimum degree three and maximum degree five.
Its rank and cycle count remain unchanged.
Its underlying simple graph is a minor of $G$, and
\[
 |V(\Gamma)|\le2r-2.
\]

Put $k=\lceil\log_2r\rceil$, and choose $k$ edges $D$ outside a spanning tree of $\Gamma$.
Let $P_0$ be their endpoints, so $|P_0|\le2k$.
Deleting a vertex of degree at most five decreases rank by at most four.
Hence the graph $J=\Gamma-P_0$ has
\[
 \beta(J)\ge r-8k.
\]
Apply Lemma~\ref{lem:path-small-components} to $J$, with $K=4096h^3$.
There is a set $S$ of size at most $|V(J)|/16$ such that every component of $J-S$ has at most $K$ vertices.
Another rank estimate gives
\[
 \beta(J-S)\ge r-8k-4|S|
 \ge r/2-8k\ge r/4.
\]
The last inequality uses $r\ge32\lceil\log_2r\rceil$, which holds for $r\ge512$.

Every component of $J-S$ has rank at most $2K$.
Choose one cycle from each cyclic component.
This yields $q\ge r/(8K)$ vertex-disjoint cycles, each using at most $K$ vertices.
Each chosen vertex set has edge boundary at most $3K$ in $\Gamma$.

Now add one fresh two-edge path parallel to each selected edge in $D$.
Let $H$ be the resulting graph, and let $W$ be the union of the resulting triangles.
The added edges only meet $\Gamma$ at vertices in $P_0$.
Thus the selected cyclic regions retain their boundary bound in $H$ and avoid $V(W)$.
The small-cut estimate gives
\[
 p(H,W)\le\frac12+\frac{(3K-1)2^{3K-2}}{q}
 \le\frac12+\frac{8K(3K-1)2^{3K-2}}{r}.
\]
The triangle-flip injection from Proposition~\ref{prop:planar-cycle-error} bounds cycles meeting $D$ by $p(H,W)2^r$.
At most $2^{r-k}\le2^r/r$ cycles avoid $D$.
Together these estimates give the stated constant $C_h$.

If $c(G)\le2^{r-1}$, the result holds immediately.
For $1\le r<512$, use $c(G)<2^r$ and $C_h\ge256$.
Finally, every proper minor-closed class excludes some $K_h$ with fixed $h\ge4$.
\end{proof}

This result is a consequence of the present counting and small-cut arguments together with the classical separator theorem.
Its publication priority has not been checked.
